\documentclass[border=5pt]{standalone}
\usepackage{pgfplots}
\usepackage{pgf-pie} % 饼图：AI 代码可直接使用 \pie
\pgfplotsset{compat=1.18}
\usepgfplotslibrary{fillbetween}  % 面积图 / 堆叠面积图 / \closedcycle 填充路径
\usepgfplotslibrary{errorbars}     % \addplot +[error bars/.cd, ...] 误差棒子键
\usepackage{amsmath}
\usepackage{amssymb}
\usepackage{xcolor}

% 支持中文
\usepackage{fontspec}
\usepackage{xeCJK}
\setCJKmainfont{SimSun}
\setmainfont{Times New Roman}

\begin{document}

\begin{tikzpicture}
\begin{axis}[
    width=14cm, height=8cm,
    title={2015—2024 年中国建成区绿化覆盖率变化},
    xlabel={年份},
    ylabel={建成区绿化覆盖率（\%）},
    xmin=2014.6, xmax=2024.6,
    ymin=38, ymax=49,
    xtick={2015,2016,2017,2018,2019,2020,2021,2022,2023,2024},
    ytick={38,40,42,44,46,48},
    ymajorgrids=true,
    grid style={dashed, gray!30},
    tick align=outside,
    legend style={at={(1.03,0.5)}, anchor=west, draw=black, fill=white},
]

% 面积主体：折线下方填充淡蓝色区域
\addplot[fill=blue!15, draw=blue!70!black, thick] coordinates {
    (2015,40.1) (2016,40.8) (2017,41.5) (2018,42.3) (2019,43.0)
    (2020,43.8) (2021,44.5) (2022,45.2) (2023,46.0) (2024,46.8)
} \closedcycle;
\addlegendentry{建成区绿化覆盖率}

% 数值标注：单系列 10 个点,按交错的上下方向标注（奇数年标上方,偶数年标下方）
\addplot[only marks, mark=*, mark size=2pt, blue!70!black, forget plot,
    nodes near coords,
    every node near coord/.append style={font=\scriptsize, fill=none, draw=none, inner sep=1pt, anchor=south}
] coordinates {
    (2015,40.1) (2017,41.5) (2019,43.0) (2021,44.5) (2023,46.0)
};

\addplot[only marks, mark=*, mark size=2pt, blue!70!black, forget plot,
    nodes near coords,
    every node near coord/.append style={font=\scriptsize, fill=none, draw=none, inner sep=1pt, anchor=north}
] coordinates {
    (2016,40.8) (2018,42.3) (2020,43.8) (2022,45.2) (2024,46.8)
};

\node[anchor=north west, font=\scriptsize] at (axis description cs:0.0,-0.15)
    {数据来源：住房和城乡建设部《城乡建设统计年鉴》（2015—2024 年）};
\end{axis}
\end{tikzpicture}

\end{document}